\chapter{The PHYSICS of TIME TRAVEL}\label{ch:time}
\markboth{The Physics of Time Travel}{The Physics of Time Travel}

\chapquote{``Originally, the burden of proof was on physicists to prove that time travel was possible. Now the burden of proof is on them to prove there must be a law forbidding time travel.''}{Dr.~Michio Kaku}

% ---------------------------------------------------------------
\section{The Physics of Time Travel}\label{sec:physicstimetravel}

Time travel into the future is not controversial. You are doing it now, at one second per
second, and if you would like to do it faster you need only move quickly or sit somewhere
with a stronger gravitational field. This is not speculation. It is measured daily, and
your satellite navigation would drift by several kilometres a day if the correction were
not applied.

Special relativity gives you the first mechanism. Time passes more slowly for a moving
observer relative to a stationary one; at 87\% of light speed, clocks run at half rate. A
crew travelling at that velocity for what they experience as ten years returns to find
twenty years elapsed at home. \hi{That is forward time travel, it requires no exotic physics,
and its only obstacle is energy.} General relativity gives you the second: clocks run slower
in stronger gravity. Cosmonaut Sergei Krikalev, having accumulated 803 days in orbit, is
the most time-displaced human being alive --- by roughly 0.02 seconds. Do not spend it all
at once.

The past is a completely different problem, and here is the crucial distinction that most
popular accounts blur: travel to the past is not obviously forbidden by the equations, and
that is precisely what worries physicists.

General relativity permits solutions containing \term{closed timelike curves} --- paths
through spacetime that return to their own starting point in time as well as space. Kurt
G\"odel found one in 1949, in a rotating universe, and presented it to Einstein as a
birthday present, which tells you something about mathematicians. Others followed: the
interior of a rotating black hole; the infinite rotating cylinder of van Stockum and
Tipler; the two cosmic strings of Gott; and the wormhole of Morris, Thorne and Yurtsever,
which remains the most-discussed because it is the least absurd.

\figwrap[l]{0.38}{Embedding diagram of a traversable wormhole: two mouths joined by a throat that exotic matter would have to hold open.}{https://commons.wikimedia.org/wiki/File:LorentzianWormhole.jpg}{fig:wormhole}
A traversable wormhole is a topological shortcut --- two mouths connected by a throat (see Figure~\ref{fig:wormhole}). Move
one mouth at relativistic speed and it becomes desynchronised from the other by ordinary
time dilation. Step in one end and out the other and you have moved in time. The
difficulty is that keeping such a throat open requires \term{exotic matter} with negative
energy density, in quantities that make ``vast'' an understatement, and while negative
energy density does exist in tiny amounts as a quantum effect --- the Casimir effect is
measured --- nobody has any idea how to accumulate it.

\begin{funfact}
The Alcubierre drive, endlessly cited in this field as ``NASA's warp drive'', is a genuine
1994 paper by a genuine physicist, and its point was arguably the opposite of what
enthusiasts take from it. Miguel Alcubierre showed you can write down a metric that
contracts space ahead of a ship and expands it behind, moving the ship without local
acceleration. He also noted that it requires negative energy, that the original
formulation needed a mass-energy budget comparable to the visible universe, and that the
interior would be causally disconnected from the exterior --- meaning you could not steer.
Subsequent work has reduced the energy requirement substantially, from ``impossible'' to
``still impossible''. The paper is a demonstration that general relativity is permissive,
not a blueprint.
\end{funfact}

Then there are the objections, and they are serious. The \term{grandfather paradox} is the
famous one, and physics has produced two respectable answers. The Novikov self-consistency
principle holds that only self-consistent histories have non-zero probability --- you can
go back, but you will find that you always were part of the past, and every attempt to
create a contradiction fails for boring local reasons. The alternative, which we will meet
in Chapter~\ref{ch:dimensions}, is that the timeline branches, in which case you have not
changed your own past at all; you have merely visited someone else's.

Hawking's chronology protection conjecture is the sharper objection. It proposes that
quantum effects always intervene to destroy a closed timelike curve at the moment of its
formation --- vacuum fluctuations circulate around the loop, amplifying without limit, and
the energy density blows up and destroys the throat. It is a conjecture, not a theorem,
because proving it requires a quantum theory of gravity that we do not have. Hawking's own
empirical argument was characteristically brisk: we are not overrun by tourists from the
future.

I include this chapter, and put it early, because time travel is the most common exotic
explanation offered in Ufology after extraterrestrial visitation --- the ``ultraterrestrial''
or future-human hypothesis. It has one genuine advantage over the interstellar hypothesis:
it solves the distance problem, since a time traveller need not cross space at all. It has
one crippling disadvantage: every mechanism we can write down requires physics we cannot
perform, and several of them may be forbidden outright. Keep both facts in mind as you read
the cases that follow. Some of them are interesting. None of them require this chapter to
be true.

% ---------------------------------------------------------------
\section{Rendlesham Forest}\label{sec:rendlesham}

Over three nights in late December 1980, in a pine forest in Suffolk, England, between two
active NATO air bases holding nuclear weapons, a series of events occurred that produced
the best-documented military UFO case in European history. It is frequently called
Britain's Roswell, which is unfair to Rendlesham --- the documentation here is far better.

The core facts are not in dispute. In the early hours of 26 December, security personnel at
the east gate of RAF Woodbridge observed lights descending into Rendlesham Forest. Three
men --- Staff Sergeant Jim Penniston, Airman First Class John Burroughs, and Airman Edward
Cabansag --- were sent to investigate on foot, initially believing an aircraft had crashed.
Two nights later, on 28 December, the deputy base commander of RAF Bentwaters, Lieutenant
Colonel Charles Halt, led a team into the forest with a Geiger counter and a cassette
recorder. The recording exists. So does the memorandum Halt wrote on 13 January 1981,
titled ``Unexplained Lights'', which was released under the US Freedom of Information Act
in 1983.

The Halt memo is the case's spine, so read what it actually says. It describes a glowing
object, metallic in appearance, triangular, roughly two to three metres across the base,
illuminating the forest with a white light and having a red pulsing light on top and a bank
of blue lights underneath. It describes the object manoeuvring through the trees and
departing. It describes three depressions found the following day in a triangular pattern,
and beta/gamma readings at those depressions of 0.07 milliroentgens per hour, noted as ten
times background. It describes, on the third night, a red sun-like object which appeared to
throw off molten metal and then split into five white objects.

That is an extraordinary document written by a serving lieutenant colonel, and its
authenticity is not in question. Now the conventional case, which is also strong.

\figwrap[r]{0.42}{Orford Ness lighthouse, about eight kilometres from Rendlesham Forest, sweeping the treeline at a low angle.}{https://commons.wikimedia.org/wiki/File:Orfordness_Lighthouse.jpg}{fig:orfordness}
The forest lies about eight kilometres from Orford Ness lighthouse, whose beam sweeps at a low angle across exactly that terrain (see Figure~\ref{fig:orfordness}), and whose flash rate matches the pulsing witnesses
described. Halt's own recording captures the team discussing the lighthouse and then
dismissing it. Astronomer Ian Ridpath established that the brilliant object low in the sky
on the third night was the star Sirius, and the ``molten metal'' and splitting behaviour is
consistent with atmospheric scintillation and with the other objects reported --- which
were at bearings matching Deneb, Vega and other bright stars. The radiation readings, on
examination of the instrument's scale, fall within the normal background range for that
area; the ``ten times background'' assessment came from an operator unfamiliar with the
instrument. The three depressions were identified by a local forester as rabbit diggings,
and there was a well-known local phenomenon of a farmer's tractor and of poachers in the
forest that night.

\begin{funfact}
The Ministry of Defence's released files show that the case was assessed as being of no
defence significance --- and also show that nobody in the MoD conducted anything resembling
an investigation before reaching that conclusion. Both the believers and the debunkers of
Rendlesham have relied on the same absence: the government neither confirmed nor examined
anything. The 2001 release of the ``Rendlesham file'' was greeted as vindication by both
sides simultaneously, which is usually a sign that the file contains very little.
\end{funfact}

Then there is the part that belongs in this chapter. Jim Penniston has stated --- not in
1980, but in accounts developed from the 1990s onward, some under hypnosis --- that he
touched the craft, that its surface was warm and covered with symbols resembling Egyptian
hieroglyphs, and that he received a download of binary code. That claim, and what came of
it, is the subject of the next three sections.

My position on Rendlesham: something happened in that forest, several trained men were
frightened by it, and Halt's contemporaneous recording is the closest thing this field has
to an honest primary document. The lighthouse and stellar explanations account for the
majority of the reported observations and cannot be waved away. The residue --- and there is
a residue, chiefly the close-range 26 December encounter --- is contaminated beyond repair
by thirty years of hypnosis, book deals, and mutual contamination between witnesses who
have been in the same room at the same conferences for decades. This is the tragedy of the
case, and Section~\ref{sec:investigativetechniques} told you exactly how it happened.

% ---------------------------------------------------------------
\section{Understanding Binary}\label{sec:understandingbinary}

To evaluate what comes next you need to understand binary, and it takes about two minutes,
so let us do it properly.

Our ordinary numbers are base ten: each position is worth ten times the one to its right,
because we have ten fingers and no better reason. Binary is base two: each position is
worth twice the one to its right, and only the digits 0 and 1 exist. The place values are
1, 2, 4, 8, 16, 32, 64, 128, and so on. To read a binary number, add up the place values
wherever there is a 1.

\begin{center}\small
\setlength{\tabcolsep}{1.1em}
\begin{tabular}{@{}ll@{}}
\toprule
\textbf{Binary} & \textbf{Decimal}\\
\midrule
\texttt{0000\,0001} & 1\\
\texttt{0000\,1010} & 8 + 2 = 10\\
\texttt{0100\,0001} & 64 + 1 = 65\\
\texttt{0110\,0101} & 64 + 32 + 4 + 1 = 101\\
\bottomrule
\end{tabular}
\end{center}

Computers use binary because a switch is either on or off, and a wire either carries
current or does not. That is the entire reason. There is nothing mystical about it, and
nothing universal about it either, which will matter shortly.

To store text, you need a convention mapping numbers to characters. The dominant one is
\term{ASCII}, agreed in 1963 by an American standards committee. In ASCII, capital A is 65,
lowercase a is 97, and the digit 1 is 49. Those values are not derived from anything. They
are the result of committee decisions about teleprinter compatibility in the early 1960s.
Had the committee met on a different day, A might be 71.

Here is the part of the table that matters for what follows. If you want to check any claim in
this chapter yourself --- and you should, because this is one of the few places in Ufology where
checking is trivially possible --- you can decode a purported alien message by hand with nothing
but this page and a pencil.

\begin{center}\small
\setlength{\tabcolsep}{0.45em}
\begin{tabular}{@{}c r l @{\hspace{1.5em}} c r l @{\hspace{1.5em}} c r l@{}}
\toprule
\textbf{Ch} & \textbf{Dec} & \textbf{Binary} &
\textbf{Ch} & \textbf{Dec} & \textbf{Binary} &
\textbf{Ch} & \textbf{Dec} & \textbf{Binary}\\
\midrule
A & 65 & \texttt{0100\,0001} & J & 74 & \texttt{0100\,1010} & S & 83 & \texttt{0101\,0011}\\
B & 66 & \texttt{0100\,0010} & K & 75 & \texttt{0100\,1011} & T & 84 & \texttt{0101\,0100}\\
C & 67 & \texttt{0100\,0011} & L & 76 & \texttt{0100\,1100} & U & 85 & \texttt{0101\,0101}\\
D & 68 & \texttt{0100\,0100} & M & 77 & \texttt{0100\,1101} & V & 86 & \texttt{0101\,0110}\\
E & 69 & \texttt{0100\,0101} & N & 78 & \texttt{0100\,1110} & W & 87 & \texttt{0101\,0111}\\
F & 70 & \texttt{0100\,0110} & O & 79 & \texttt{0100\,1111} & X & 88 & \texttt{0101\,1000}\\
G & 71 & \texttt{0100\,0111} & P & 80 & \texttt{0101\,0000} & Y & 89 & \texttt{0101\,1001}\\
H & 72 & \texttt{0100\,1000} & Q & 81 & \texttt{0101\,0001} & Z & 90 & \texttt{0101\,1010}\\
I & 73 & \texttt{0100\,1001} & R & 82 & \texttt{0101\,0010} & \textvisiblespace & 32 & \texttt{0010\,0000}\\
\bottomrule
\end{tabular}
\end{center}

Two shortcuts are worth memorising. Lowercase letters sit exactly 32 above their capitals, which
in binary means flipping a single bit --- \texttt{0100\,0001} is A and \texttt{0110\,0001} is a. The
digits 0 to 9 run from 48 to 57, so the character \texttt{1} is \texttt{0011\,0001} and not
\texttt{0000\,0001}; the numeral you can see and the number the machine counts with are different
things, and conflating them has produced more than one confident decoding of a message that says
nothing.

One more thing, and it is the thing almost nobody in this argument has bothered to work
out: decoding and encoding are not the same job. \hi{A decoder can be operated by anyone;
an encoder has to be memorised.} To read a stream of ones and zeros you need nothing but the
table above, or in the modern era a free web page. To \emph{write} one --- to sit down with a
blank notebook and produce the digits for the word EXPLORATION --- you must know, without
looking anything up, that E is 69 and X is 88 and P is 80, and you must then convert each of
those to eight binary digits in the correct order, eleven times over, without losing your
place. That single word is 88 digits long:

\begin{center}\footnotesize\ttfamily
0100\,0101\ \ 0101\,1000\ \ 0101\,0000\ \ 0100\,1100\ \ 0100\,1111\ \ 0101\,0010\ldots
\end{center}

\noindent and it is one word out of roughly fifty.

The last thing you need is a sense of what damage looks like, because damaged binary is about
to become the whole argument. Errors in a bit stream are not all alike, and the two kinds leave
different fingerprints. A \term{bit error} flips one or two digits inside an otherwise correct
character, and the result is a wrong letter that is numerically close to the right one. A
\term{framing error} is worse and more revealing: if a whole character is dropped, or a group
boundary is misplaced by a few digits, everything downstream of the slip decodes into garbage
until the reader re-synchronises.

\begin{center}\small
\setlength{\tabcolsep}{0.8em}
\begin{tabular}{@{}l l l l@{}}
\toprule
\textbf{Intended} & \textbf{Binary} & \textbf{Received as} & \textbf{Digits wrong}\\
\midrule
O (79) & \texttt{0100\,1111} & C (67) \texttt{0100\,0011} & 2\\
S (83) & \texttt{0101\,0011} & T (84) \texttt{0101\,0100} & 3\\
R (82) & \texttt{0101\,0010} & \emph{nothing} & 8 (character lost)\\
\bottomrule
\end{tabular}
\end{center}

Hold on to that table. A forger garbling his own work damages \emph{letters}, because letters
are what he is thinking in. A stream that has genuinely been mis-copied is damaged at the level
of \emph{digits and groups}, because digits are what was actually on the page. The difference is
visible, it is countable, and in the next section it is the only test available to us that does
not depend on trusting anybody.

\begin{funfact}
This is the point that demolishes an entire genre of claim. If a message from a
non-human intelligence arrived encoded in eight-bit ASCII, the correct conclusion would not
be ``they can speak English''. It would be ``a human being wrote this''. ASCII is a
parochial artefact of Cold War American office equipment. Expecting an alien or a future
civilisation to use it is like expecting them to file their reports in triplicate. There is
more on genuinely universal encoding in Section~\ref{sec:mathlang}, and the answer is
mathematics --- primes, physical constants, ratios --- not character sets.
\end{funfact}

Binary itself is a reasonable candidate for a shared representation, since base two is the
minimal positional system and any species that builds switches will find it. Binary
\emph{plus ASCII} is a fingerprint. Keep that distinction sharp for the next section.

% ---------------------------------------------------------------
\section{The Impossible Notebook}\label{sec:impossiblenotebook}

Jim Penniston kept a notebook. He has stated that in the hours after the 26 December
encounter he sketched the object he had touched, copied the symbols from its surface, and
then --- over the following days --- compulsively wrote out pages of ones and zeros, without
knowing why or what they meant. The notebook pages have been photographed and published,
and the sketch of the craft with its glyphs is one of the more striking artefacts in the
field.

Sixteen of those pages are binary, and it is worth putting the volume of data in terms a
modern reader can feel. The message that was deciphered in 2010 in Sgt.\ Jim Penniston's
notebook was exactly 16 pages long and roughly equates to 302 individual characters (yes
spaces included). Let's just say an even 300. \hi{The message at 300 letters long would
require 2400 individual 1's or 0's to comprise the data within and would take up 2.4 KB of
space on your computer hard drive.} (Because each 1 or 0 represented in say a text file
requires 8 bits in itself to record). That's over half the entire storage capacity of NASA's
state of the art guidance system for the Gemini Missions conducted just 15 years prior to the
supposed interaction with a UFO in Rendlesham Forest! All of it produced over several days by
a security policeman with no programming background, no ASCII table in front of him, and, by
his own account, no idea what he was writing.

In 2010, a researcher named Nick Ciske ran the digits through an ASCII decoder. They
produced English text. Penniston later published his own composite reading of all sixteen
pages, and --- to his considerable credit, because almost nobody in this field does this ---
he colour-coded it, marking which characters fell straight out of the binary, which were
damaged, which were unreadable, and which were his own guesses. That chart is the single
most useful document in the case, and it is reproduced below, set vertically so that each
line can be inspected on its own.

\subsection*{The transmission, line by line, with the faults marked}

\newcommand{\pnd}[1]{{\color{inkblue}#1}}%
\newcommand{\pnq}[1]{{\color{gapred}\bfseries #1}}%
\newcommand{\pnu}{{\color{gapred}\bfseries ???}}%
\newcommand{\pni}[1]{{\color{ufogreendark}\bfseries #1}}%
\newcommand{\pnp}{{\color{amber}\bfseries .}}%

\begin{center}\footnotesize
\setlength{\tabcolsep}{0.5em}
\renewcommand{\arraystretch}{1.18}
\begin{tabular}{@{}r@{\hspace{0.7em}}l l@{}}
\toprule
& \textbf{Decoded line} & \textbf{Penniston's identification}\\
\midrule
\color{midgrey}1 & \ttfamily\pnd{EXPLORATION} \pnq{C}\pni{F} \pnd{HUMANITY} \pnd{666} \pnd{8100} & \\
\color{midgrey}2 & \ttfamily\pnd{52}\pnp\pni{09}\pnd{42532} \pnd{N} \pnd{13}\pnp\pnd{131269} \pnd{W} & Hy Brasil\\
\color{midgrey}3 & \ttfamily\pnd{CONTINUOUS} \pnq{F}\pnd{CR} \pnd{PLANETARY} \pnd{ADVAN}\pnu & \\
\color{midgrey}4 & \ttfamily\pnd{F OUR COODINATE CONTINU OT} \pnq{UQS} \pnd{C}\pnq{b}\pnd{PR BEFORE} & \\
\color{midgrey}5 & \ttfamily\pnd{16}\pnp\pnd{763177} \pnd{N} \pnd{89}\pnp\pnd{117768} \pnd{W} & Caracol, Belize\\
\color{midgrey}6 & \ttfamily\pnd{34}\pnp\pnd{80027}\pni{2} \pnd{N} \pnd{111}\pnp\pni{8435}\pnd{67} \pnd{W} & Sedona, Arizona\\
\color{midgrey}7 & \ttfamily\pnd{29}\pnp\pnd{977836} \pnd{N} \pnd{31}\pnp\pnd{131649} \pnd{E} & Great Pyramid, Giza\\
\color{midgrey}8 & \ttfamily\pnd{14}\pnp\pnd{701505} \pnd{S} \pnd{75}\pnp\pni{167}\pnd{043} \pnd{W} & Nazca Lines, Peru\\
\color{midgrey}9 & \ttfamily\pnd{36}\pnp\pnd{256845} \pnd{N} \pnd{117}\pnp\pnd{100632} \pnd{E} & Tai Shan, China\\
\color{midgrey}10 & \ttfamily\pnd{37}\pnp\pnd{110195} \pnd{N} \pnd{25}\pnp\pnd{372281} \pnd{E} & Portara, Naxos, Greece\\
\color{midgrey}11 & \ttfamily\pnd{EYES OF YOUR} \pni{EYES} & \\
\color{midgrey}12 & \ttfamily\pnd{ORIGIN} \pnd{52}\pnp\pni{0}\pnd{942532} \pnd{N} \pnd{13}\pnp\pnd{131269} \pni{W} & Hy Brasil\\
\color{midgrey}13 & \ttfamily\pnd{ORIGIN YEAR} \pnd{8100} & \\
\bottomrule
\end{tabular}
\end{center}

\begin{center}\footnotesize
\begin{minipage}{0.86\linewidth}
{\color{ufogreenink}\scshape Key.}\ 
\pnd{Blue} --- decoded directly from the binary, no judgement required.\ 
\pnq{Red} --- questionable characters, in stretches where the stream is faulty.\ 
\pnu\ --- characters that could not be read at all.\ 
\pni{Green} --- Penniston's own interpretation, supplied where the data are damaged.\ 
\pnp\ --- decimal points inserted by hand; the stream contains no punctuation.\ 
Black --- his commentary, not part of the message.
\end{minipage}
\end{center}

Read the display rather than the summary and three things become obvious that the usual
quotation hides. First, the faults are not evenly spread: lines 1, 3 and 4 carry almost all
of the damage, and the ten coordinate lines are very nearly clean. Second, the decimal points
in every coordinate were \emph{put there by a human being}; the binary contains no punctuation
at all, so the difference between 52.0942532 and 520.942532 is Penniston's decision, not the
message's. Third, and this is the number the whole case turns on: count the marked characters
against the total and \hi{the stream decodes to intelligible English at roughly 95 per cent
fidelity}. About one character in twenty is damaged. Nineteen in twenty are not.

That figure is the actual mystery here, and it is not the mystery either side thinks it is.

The coordinates have been matched to various locations of significance --- Hy Brasil, a
legendary island west of Ireland; the Great Pyramid; Sedona; the Nazca lines. The year 8100
is taken to indicate that the visitors were future humans, which is why this case sits in a
chapter about time travel rather than one about extraterrestrials.

Now. Section~\ref{sec:burdenproof} says the burden of proof falls on the claimant, and Section~\ref{sec:doubtdrivenskepticism} says apply
your hardest scrutiny to the evidence you like best. So:

The message is in English. Not merely encoded in an American character set --- actually
composed in English, with English word order and English spelling. It uses decimal degrees
of latitude and longitude, a coordinate system defined by a British prime meridian and a
Babylonian sexagesimal circle, and it specifies a year in a calendar reckoned from a
Christian event. A message from 8100~AD, or from anywhere else, that arrives in
English-language ASCII with WGS-84 coordinates and an Anno Domini date has told you
precisely one thing about its origin, and it is not a flattering thing.

The provenance is worse. The notebook was not made public in 1980, or 1985, or 1990. The
binary pages entered the record decades later, after Penniston had undergone hypnotic
regression --- see Section~\ref{sec:hypnosis} --- and after he had become a regular figure
on the conference circuit. The decoding was performed by someone who knew what case he was
working on and what answer would be interesting. There is no chain of custody, no dated
third-party examination of the pages, and no independent verification of when the digits
were written. Section~\ref{sec:deconstructingpeerreview} covered exactly this: an analysis performed by a partisan on an
unauthenticated artefact is not evidence, however competent the analyst.

And the decoding itself is not clean. The digit sequences require judgement about where
characters begin and end, and different investigators reading the same pages have produced
different text. When a decoding requires the decoder to make choices, and the decoder knows
what he hopes to find, Section~\ref{sec:confirmationbias} tells you what the output will be.
The display above is the honest form of that admission, and it is why I have printed it in
full rather than the tidy one-line quotation.

\subsection*{The question everybody is asking is the wrong one}

Having said all that, the standard skeptical treatment of this case --- mine included, until I
sat down and actually counted --- stops one question too early. Everyone argues about whether
the message is a genuine transmission. Nobody asks how the pages got written.

\actualq[question]{Did Jim Penniston receive a binary message from a craft in Rendlesham
Forest?}{How did a man with no programming background, no ASCII table, and no reason to own
one put roughly 2,400 binary digits on paper that decode to intelligible English at around 95
per cent fidelity?}

That second question has to be answered by everybody, whichever side they are on, and it is
harder than it looks in both directions.

Start with what the task actually demands, because the previous section set it up. Decoding is
free; encoding is not. A person writing this material from scratch is not typing English --- he
is performing a two-stage conversion, letter to decimal to eight binary digits, several hundred
times in a row, on paper, with no error correction and no way to check his work. \hi{Nobody
writes 2,400 correct binary digits by accident, and nobody writes them from memory either.}
The world record for memorising binary digits, set under competition conditions by trained
memory athletes using the method of loci, stands in the thousands --- but those are people who
have spent years building the technique, working from a random string they are handed, and
reciting it back within the hour. That is not what is claimed here. Penniston claims the digits
arrived unbidden and came out over days.

Then there is the clock. The handwriting speed of the average human is approximately 68
characters per minute. Given that information, one could deduce that writing numbers
frivolously into a notebook without thought as Penniston claims he did is a task that would
take him approximately 35 minutes to accomplish. Imagine that! Writing nothing but 1's and
0's at random onto the pages of a small notebook for nearly a solid 40 minutes! I think given
the accuracy of the translated text to convert to legible English indicates it's obvious that
such a task as writing over two thousand 1's and 0's was NOT done randomly!

So there are three possibilities, and they should be laid out plainly rather than assumed away.

\begin{enumerate}\setlength{\itemsep}{2pt}
\item \textbf{He copied them from somewhere.} Somebody --- himself, a collaborator, or a source
he has never named --- composed the English sentences first, ran them through an encoder, and
produced the digit stream, which was then transcribed into the notebook by hand. This requires
no unusual ability at all. It requires an encoder, a few hours, and a willingness to write.
\item \textbf{He wrote them from memory of something seen.} This is his own account, and on the
numbers it is the least plausible of the three, because nothing in the human memory literature
supports the untrained recall of thousands of undifferentiated binary digits days after a
thirty-second exposure.
\item \textbf{The digits were generated by a process nobody has described.} Left on the table
for completeness, but an unspecified mechanism is not an explanation, and
Chapter~\ref{ch:language} says so.
\end{enumerate}

Option one is dull, and it is also the only one that survives contact with the fidelity figure.
Here is why the 95 per cent is the tell rather than the marvel. Ask what damaged data should
look like. A stream genuinely received through a noisy channel, or genuinely mis-copied from a
half-remembered image, degrades at the level of digits and groups: one dropped digit throws the
frame, and everything after it is rubbish until the reader happens to re-synchronise. Whole
sentences do not survive that. Yet the flaws in the display above are almost all
\emph{single letters} sitting inside otherwise perfect words --- a C where an O belongs, an
intact CONTINUOUS followed by a garbled fragment, an unreadable run at the end of ADVAN and
then ten coordinate lines that are essentially flawless. That is the signature of a transcription
error made by a hand copying digits it does not understand, not of a corrupted channel. The
coordinates come out cleanest of all, which is precisely backwards for a damaged transmission
and precisely right for a document whose author cared most about the coordinates.

And the fidelity is too high to be innocent in the other direction as well. Random noise does
not decode to English at all; eight-bit ASCII is a sparse code, and the overwhelming majority
of byte values are not printable letters. Getting nineteen characters in twenty to land inside
the alphabet is not something a stream can do by luck. \hi{Whatever produced these pages had
an ASCII table in the loop.} The only live question is whether it was in the loop in 1980, in
1994, or in 2010.

\subsection*{Hy Brasil, and the odds people keep quoting wrongly}

The strongest-sounding item in the whole display is the repeated coordinate, given twice ---
once plainly, once as ORIGIN --- and landing near the position associated with Hy Brasil, the
phantom island charted west of Ireland from the fourteenth century onward and now generally
identified with the Porcupine Bank, a submerged plateau roughly 200 kilometres off the Irish
coast. Believers present this as a coincidence too large to accept. Skeptics reply that any
coordinate hits something legendary, which the funfact below addresses.

Both are sloppy, so here is the honest arithmetic. If the coordinate had been drawn at random,
the chance of landing within a hundred kilometres of any \emph{particular} named site is on the
order of one in ten thousand. That is the number believers quote. But nobody specified Hy Brasil
in advance. The coordinate was decoded first and the site was found afterwards, from a list of
candidate anomalies that is effectively unbounded --- every megalith, every phantom island,
every ley intersection, every crash site. Once you allow the target to be chosen after the shot,
the odds collapse, and \hi{a match found after the fact is not a prediction, it is a search
result.} That is Section~\ref{sec:correlationvscausation} in one line.

What is genuinely interesting about the Hy Brasil line has nothing to do with probability. It is
that the coordinate needed a decimal point, and the decimal point is Penniston's. Slide it one
place and the position leaves the Atlantic entirely. A message that cannot locate itself without
human punctuation is not a set of coordinates; it is a string of digits that a human being turned
into coordinates.

\subsection*{8100: the one part of the message that behaves like a message}

The last two lines are the part I cannot dismiss as cleanly as I would like, and honesty
requires saying so. Set out the structure rather than the content:

\begin{center}\small
\setlength{\tabcolsep}{0.9em}
\begin{tabular}{@{}l l l@{}}
\toprule
\textbf{Line} & \textbf{Content} & \textbf{What it specifies}\\
\midrule
ORIGIN & 52.09 N, 13.13 W & a place\\
ORIGIN YEAR & 8100 & a time\\
\bottomrule
\end{tabular}
\end{center}

Taken together those two lines do not give a location and they do not give a date. They give
\hi{a single point in space and time}, which is the only form of address that means anything to
a traveller for whom the two are one quantity. That is a coherent thing to send. It is also,
note, exactly the shape of message that somebody who had decided in advance that the visitors
were future humans would compose --- and 8100 appears twice, once at the head of the stream
(where it is usually misquoted as 8145) and once at the foot, which is how a person writing a
document frames it, with the headline restated at the end.

I do not know which of those it is. The evidence in the pages themselves cannot settle it,
because the pages have no provenance. What the pages \emph{can} settle is the fidelity question,
and that answer points one way.

\begin{funfact}
Between them the decoded coordinate lines have been matched to Hy Brasil, Caracol, Sedona,
Giza, Nazca, Tai Shan and Naxos --- seven sites of anomalous interest on four continents, all
found after the digits were read rather than before. That is a statistical inevitability rather
than a mystery. Pick any latitude and longitude on Earth and something legendary lies within a
few hundred kilometres, because legends are dense on the ground. This is the
Section~\ref{sec:correlationvscausation} problem wearing a compass.
\end{funfact}

I regard the notebook as the single most instructive artefact in this book, and not because
it is genuine. It is instructive because it is \emph{almost} good enough to be compelling,
and because it fails on a point --- the character encoding --- that is checkable by anybody
who has read the previous section. That is what a real methodology gets you. You do not
need access to Suffolk, or to Penniston, or to a laboratory. You need base two and about
five minutes.

% ---------------------------------------------------------------
\section{John Titor}\label{sec:johntitor}

Beginning in November 2000, a poster on internet forums including the Time Travel Institute
and Art Bell's BBS claimed to be a soldier from the year 2036, sent back to 1975 to retrieve
an IBM 5100 computer needed to debug legacy UNIX systems, and stopping in 2000--2001 for
personal reasons. He posted under the name John Titor, produced diagrams of a
``C204 gravity distortion time displacement unit'' installed in a Chevrolet, photographs of
the device, and a scanned military insignia. He answered questions for several months and
then, in March 2001, said he was going home and stopped posting.

The Titor material is the most successful piece of internet folklore the UFO-adjacent world
has produced, and it is worth studying as folklore rather than as physics.

His predictions failed, and they failed specifically. He stated that a civil war would begin
in the United States around 2004--2005 over a disputed federal authority and escalate into a
brief nuclear exchange with Russia in 2015, reducing the country to a rural agrarian society
of five regional capitals. He stated that CERN would discover the basis of time travel around
2001 with the production of miniature black holes. He stated the 2004 Olympics would be the
last. None of this happened. There was no 2015 nuclear war, the Olympics continue, and CERN's
actual 2012 discovery was the Higgs boson, which is not a portal.

His defenders invoke the branching-worlds escape --- he said explicitly that his timeline
diverged from ours by a few per cent, meaning no prediction can ever falsify him. Note what
that does. \hi{An unfalsifiable claim is not a strong claim, it is an empty one}, and
Section~\ref{sec:claimsmadeslipperyslope} covered the mechanics of this retreat.

The technical details are the tell. The IBM 5100 detail is the celebrated ``hit'': Titor
said the machine had an undocumented ability to emulate older IBM mainframe code, which was
not public knowledge and was confirmed by an IBM engineer. It is a real feature. It was also
discussed in IBM technical circles and in print, and the person best placed to know it would
be a computing hobbyist in 2000 --- not a soldier in 2036, for whom a fifty-year-old
portable computer would be an artefact of no more personal familiarity than a Victorian
loom is to us.

\begin{funfact}
In 2003 an Italian researcher, working with a lawyer, traced the Titor material to a
Florida entertainment company and to two brothers, one of whom had a background in computing
and the other in law. A linguistic analysis published in 2009 by John H\aa{}kan Larsson found
that the writing style matched. No confession has ever been made, and the family has
consistently declined to comment --- which is, in fairness, exactly what innocent people also
do.
\end{funfact}

What makes Titor interesting is not the claim but the shape of it. He was modest,
non-messianic, occasionally rude, refused to profit, made no attempt to be believed, and
left on schedule. That is a far more sophisticated performance than the usual prophet
routine, and it is why the story has outlived every prediction in it. If you ever wonder why
this field is so hard to clean up, consider that a well-constructed fiction with a
consistent voice will always beat a badly written truth.

% ---------------------------------------------------------------
\section{The Bermuda Triangle}\label{sec:bermudatriangle}

A book about ufology has two obvious places to put the Bermuda Triangle --- the region bounded
roughly by Miami, Bermuda and Puerto Rico, in which ships and aircraft are said to disappear at
an anomalous rate. One is with the water cases: the USOs of Section~\ref{sec:usos}, the lost
ships, the Atlantis material. That is where most authors file it, and it is the wrong shelf. The
Triangle is in this chapter, between John Titor and a Nazi bell, for exactly one reason, and the
reason has a name: \term{electronic fog}. Somebody flew into the Triangle, came out the other
side, and reported not that he had seen something in the sea but that he had lost time.
Everything else in the legend is bookkeeping. That claim is a claim about spacetime, so it belongs
next to the physics.

The man is Bruce Gernon, and his flight is the single best reason the Triangle survived Kusche.
On 4 December 1970 Gernon took off from Andros Island in the Bahamas for Palm Beach, roughly
250~miles, in a Beechcraft Bonanza A36, with his father and a business associate aboard. Ahead
of him sat a lens-shaped cloud that appeared to be growing as he climbed. He could not get over
it at 11,500~feet. Skirting it, he found what he described as a tunnel through the cloud mass,
perhaps a mile wide at the mouth and narrowing, and flew into it. The walls contracted around
the aircraft. The compass began to rotate slowly and would not settle; every electronic
instrument on board went unreliable at once. He came out of the far end into a uniform greyish
white with no horizon, no sea and no sky --- and then out of that into clear air over Miami
Beach, having been airborne 34~minutes on a flight that should have taken well over an hour. He
was, depending on which of his accounts you read, somewhere between thirty and forty-seven
minutes early, with about a hundred miles unaccounted for and the fuel to match. Miami approach,
he says, could not see him where he said he was.

Gernon coined the phrase ``electronic fog'' for the grey substance and has spent five decades
developing it, latterly with the co-author Rob MacGregor, in \emph{The Fog} and \emph{Beyond the
Bermuda Triangle}. The proposal is that this is not ordinary cloud but something that attaches
itself to a vessel and travels with it, in the manner of St~Elmo's fire, disorienting the crew
and --- in the strong version --- opening a shortcut through space and time. That is a
falsifiable-sounding claim about the behaviour of spacetime in a specific corridor of the
Atlantic, and it is the only version of the Triangle that earns a hearing in a chapter like this
one.

So let us give it the hearing, and then do to it what has to be done to the rest.

The legend Gernon inherited was already the most commercially successful geographic mystery in
history, and as an anomaly it was already manufactured --- we can say that with unusual
confidence, because somebody did the work. The name was coined by Vincent Gaddis in a 1964
\emph{Argosy} magazine article. The concept was made famous by Charles Berlitz's 1974 bestseller
\emph{The Bermuda Triangle}, which sold in the millions. And in 1975 a librarian in Arizona named
Lawrence David Kusche did something nobody in the genre had bothered to do: he went to the primary
sources for every case Berlitz and his predecessors cited --- Lloyd's records, US Coast Guard
reports, newspaper archives, contemporary weather data.

\figwrap[l]{0.38}{The Bermuda Triangle as usually drawn --- the boundaries vary by tens of thousands of square miles between sources.}{https://commons.wikimedia.org/wiki/File:Bermuda_Triangle.png}{fig:triangle}
His findings, published as \emph{The Bermuda Triangle Mystery: Solved}, are devastating in their specificity (see Figure~\ref{fig:triangle}). Some vessels listed as vanishing in calm weather had disappeared during
documented severe storms. Some had not disappeared in the Triangle at all --- one loss placed
there occurred in the Pacific. Some had not disappeared: at least one ship reported as lost
with all hands had arrived safely in port. In several cases the ``mystery'' consisted
entirely of details Berlitz had omitted, such as a wreck having been found. Kusche's summary
is the best sentence written about this subject: the legend was ``a manufactured mystery,''
perpetuated by writers who used careless research and sometimes deliberate
sensationalism.

And the statistics settle it. Lloyd's of London stated that the Triangle is no more dangerous
than any comparable stretch of ocean, and charges no premium for it. The US Coast Guard says
the same. The area happens to be one of the most heavily trafficked shipping and aviation
regions on Earth, it contains the Gulf Stream, and it sits in the Atlantic hurricane track.
A high absolute number of losses in a high-traffic, high-weather corridor is arithmetic, not
anomaly --- the Section~\ref{sec:correlationvscausation} lesson again.

\begin{funfact}
The founding case, Flight~19 --- five TBM Avengers lost on a training flight from Fort
Lauderdale on 5 December 1945 --- is genuinely tragic and not at all mysterious once you read
the Navy's own board of investigation. The flight leader became convinced his compass was
faulty and that he was over the Florida Keys when he was in fact over the Bahamas, and led
the formation further out to sea in deteriorating weather until they ran out of fuel after
dark. Radio transcripts capture the disorientation in real time. The board initially cited
pilot error, then amended it to ``causes or reasons unknown'' after the flight leader's
mother objected --- and that amended phrase, a compassionate gesture to a grieving parent,
became the evidentiary foundation of a global mystery. The rescue aircraft that also vanished
that night was a PBM Mariner, a type nicknamed the ``flying gas tank''; a tanker crew
reported seeing an explosion and an oil slick at the appropriate time and place.
\end{funfact}

Which brings us back to Gernon, because the statistical argument does not touch him. He is not
a Berlitz entry. He was there, he is named, he has never recanted, and nobody has caught him
lying about the date, the aircraft or the passengers. The honest reply to him is not that he is a
fraud but that his aircraft was the wrong instrument for the measurement he thinks he took. A
1970 Bonanza carried no recorder; the times are recalled times. A pilot in cloud with a rotating
compass and no horizon is in the exact condition that aviation medicine calls spatial
disorientation, and the first thing it destroys is the sense of elapsed time --- which is why he
is right that this fog kills people, and right about the mechanism, and then draws the wrong
conclusion from it. The Bahamas-to-Florida route also sits under the Gulf Stream's own weather
factory, where a light aircraft climbing round a building cumulus can pick up a substantial
tailwind it has no way of measuring, and the arithmetic of a shortened flight is ground speed,
not spacetime. Fuel burn tells you how long the engine ran, not how far the aeroplane went.

That is the whole reason the Triangle is here rather than with the USOs. The exotic explanations
for it --- time slips, dimensional portals, electromagnetic vortices --- require everything in
Section~\ref{sec:physicstimetravel} to be not merely possible but routine, in one patch of the
Atlantic, for no stated reason, and to leave no trace on any of the instrumented traffic that
crosses the same water daily. Gernon's account requires one honest man in a badly instrumented
cockpit to have misjudged an interval under conditions specifically known to make people misjudge
intervals. Kusche's requires only that authors do not check things.
Section~\ref{sec:burdenproof} tells you which of those is carrying the lighter load.

I find I do not want to be glib about it, though. Of everything in this chapter, Gernon's is the
claim I would most like to be true, and he did something none of the Triangle's authors did: he
put his own name to a first-person report, described a mechanism, and invited people to go and
look for it. That is how it is supposed to work. It is not his fault that the answer, when you
follow his own mechanism carefully, is meteorology.

% ---------------------------------------------------------------
\section{Die Glocke}\label{sec:glocke}

\figwrap[l]{0.42}{Die Glocke as the account describes it: a bell-shaped shell about three metres tall, housing two counter-rotating cylinders charged with the violet liquid ``Xerum 525''. This is my own rendering of that description and nothing more --- no photograph, plan, component or measurement of the device exists to reproduce, so what the reader is looking at is a picture of a claim rather than of an object.}{}{fig:glockeschem}
\emph{Die Glocke} --- The Bell --- is a claimed Nazi secret weapon: a metallic bell-shaped
device, roughly three metres tall, containing two counter-rotating cylinders filled with a
violet mercury-like substance code-named ``Xerum 525'', which when energised produced
intense radiation, killed several of the scientists working on it, and was variously an
antigravity device, a time machine, or a means of generating a torsion field (see
Figure~\ref{fig:glockeschem}).

The entire story traces to a single source: the Polish journalist Igor Witkowski, who
stated in his 2000 book \emph{Prawda o Wunderwaffe} that he had been shown transcripts of
the interrogation of SS General Jakob Sporrenberg by an unnamed Polish intelligence officer,
was permitted to copy them by hand, and was not allowed to retain or photograph the
documents. The story reached the English-speaking world through Nick Cook's \emph{The Hunt
for Zero Point} in 2001, and from there into hundreds of documentaries.

So the evidentiary chain is: an unnamed officer showed an author documents nobody else has
seen, which the author transcribed by hand and cannot produce. That is not a chain. That is
a single link with nothing on either end. Sporrenberg was a real person, tried and executed
in Poland in 1952 for genuine atrocities, and his actual trial record contains nothing
about any such device.

\figwrap[r]{0.38}{The Wenceslas Mine at Ludwikowice, commemorated on site. The colliery was real, its wartime forced-labour history is documented, and the concrete ring nearby is a cooling tower footing --- no photograph of Die Glocke exists, because no evidence of the device does.}{https://commons.wikimedia.org/wiki/File:Ludwikowice._Tablice4.jpg}{fig:glocke}
The supporting evidence usually offered is a concrete structure near the Wenceslas Mine at
Ludwikowice in Poland (see Figure~\ref{fig:glocke}) --- a ring of eleven concrete pillars,
presented in documentaries as a test rig for the Bell. Industrial archaeologists identify it as
a cooling tower support for the adjacent mine's power plant, of a standard design found at other
sites. It looks exotic only if you have never seen a cooling tower footing.

\begin{funfact}
The Nazi wonder-weapon genre exists because the real programmes were genuinely astonishing.
The V-2 was the first object to reach space. The Me~262 was the first operational jet
fighter. The Type~XXI submarine was a generation ahead of anything afloat. The Horten Ho~229
was a flying wing with a plausibly reduced radar signature, and its shape is the ancestor of
the B-2. When a regime demonstrably built five impossible things, a sixth is an easy sell.
The difference is that the real programmes left factories, blueprints, prototypes, and
thousands of witnesses. Die Glocke left one journalist's handwriting.
\end{funfact}

There is also a moral point here, and I do not usually make them. The Nazi-antigravity genre
frequently shades into admiration --- for the engineering, the secrecy, the supposed
superiority --- and from there into the Vril and Thule occult material and eventually into
the territory I address in Section~\ref{sec:aryan}. Be aware of where this road goes. A
great deal of ``suppressed Nazi technology'' literature is doing work that has nothing to do
with physics.

% ---------------------------------------------------------------
\section{Kecksburg UFO Incident}\label{sec:kecksburg}

I have put Kecksburg immediately after Die Glocke deliberately, because the two are joined in the
popular literature and I want to separate what is joined badly from what is joined honestly.

On the evening of 9 December 1965, at about 4:43~p.m.\ Eastern, a brilliant fireball was seen over
six American states and the Canadian province of Ontario. It dropped hot metal debris that started
grass fires in Michigan and Ohio, and it was reported by thousands of people. Nobody disputes any of
that. What is disputed is what happened next in Kecksburg, an unincorporated village in Westmoreland
County, Pennsylvania, about thirty miles south-east of Pittsburgh, where witnesses said something
came down in the woods.

I am not going to spend this section arguing about whether an object landed, because that question
is unanswerable with what exists. I want to ask a narrower and much better question, and it is the
question that makes Kecksburg worth a section at all:

\pullquote{Why were men in officially badged uniforms, identifying themselves as NASA, standing
in a Pennsylvania wood that night --- and why can the agency that badged them produce no record of
it?}

\Needspace*{0.40\textheight}
\subsection*{What the witnesses said}

\figwrap[l]{0.34}{The monument in Kecksburg: a fibreglass acorn built to the witnesses' description, complete with the band of glyph-like markings. It is a sculpture of a testimony, not of a recovered object --- which is precisely the difficulty.}{https://commons.wikimedia.org/wiki/File:Kecksburg_UFO.JPG}{fig:kecksburg}
The local volunteer fire department went into the woods. James Romansky, a young fireman that
night, described an object the size of a Volkswagen, acorn- or bell-shaped, with a raised band
around the base carrying markings he compared to Egyptian hieroglyphics --- the description the
village later built in fibreglass and stood beside its fire hall (see Figure~\ref{fig:kecksburg}). He also described what
happened when the vehicles arrived: two men in unmarked clothing told the firemen to get out, in
words to the effect that they were taking command. Jim Mayes reported blue pulsing lights in the
woods. Bill Weaver described men in what he called moon suits and a box roughly five feet long being
removed. Lillian Hays's farmhouse near the site was used by the military as a base of operations,
and she later noticed that the calls made from her telephone that night never appeared on her bill.

And then there is John Hays, who was ten years old. He said he watched a flatbed truck pass his
bedroom window carrying a covered object he described as about the size of a Volkswagen Beetle. He
told that story on \emph{Unsolved Mysteries} in 1990, and he told the same story again in 2003 for
the Sci-Fi Channel documentary \emph{The New Roswell: Kecksburg Exposed}, presented by Bryant
Gumbel --- thirteen years apart, without material drift. That consistency is worth something. It is
not worth an object; a ten-year-old cannot tell you what is under a tarpaulin. What it is worth is
the recovery: something was loaded and driven away by people with the authority to close a road.

Stan Gordon, who has worked the case longer than anyone, obtained the Air Force file under FOIA.
The Air Force position is that a three-man team from Oakdale searched and found nothing. Gordon
noted that the log of the Pittsburgh-area radar squadron is missing its entry for 9 December 1965
--- the one day the log would matter.

\subsection*{What the official explanation says}

The astronomers were unanimous and quick. Paul Annear, William Bidelman, Von Del Chamberlain and
others identified the fireball as a bolide --- a meteor bright enough to fragment audibly. In 1967
Chamberlain and Krause published a triangulation in the \emph{Journal of the Royal Astronomical
Society of Canada} which gave a steep descent angle and a terminal point near Lake Erie, in the
region of Windsor, Ontario. That is a long way from Kecksburg.

The rival hard-hardware candidate was Kosmos~96, a failed Soviet Venus probe. NASA's own National
Space Science Data Center now states that this cannot be right: the Kosmos~96 orbit decayed hours
earlier over Canada, and the Kecksburg object's trajectory was too steep to match a decaying
satellite. MUFON investigators John Ventre and Owen Eichler have argued for a General Electric
Mark~2 re-entry vehicle from a failed American reconnaissance satellite, which would explain both
the acorn shape and the urgency of the recovery rather elegantly, and which remains unproven.

So the physical evidence points at a bolide, and the recovery testimony points at somebody
collecting something. Those two facts are not actually in conflict --- a hot fragment on the ground
is worth collecting, and the men who collect it are not obliged to explain themselves.

\subsection*{The men who said they were NASA}

Here is where honesty costs the story something. James Oberg, writing in 2008, made the argument
that has to be answered: in the 1960s, personnel recovering space debris routinely identified
themselves as NASA whether or not they were, and the men at Kecksburg were most probably Air Force
in civilian clothes. That is not a dismissal invented to be convenient. It is how the period
actually worked, and it explains the badge without requiring anything exotic.

What it does not explain is the paperwork. In December 2005 the NASA spokesman David Steitz told
the Associated Press that the Kecksburg object had been a Russian satellite, and that the relevant
records had been lost. That is a strange thing for a spokesman to say --- it concedes a
determination while disclaiming the file that supports it. Journalist Leslie Kean, backed by the
Sci-Fi Channel, had already sued under FOIA in 2003. Judge Emmet Sullivan described NASA's handling
of the matter as a ball of yarn and remarked that heads should roll. The agency settled in October
2007 and completed its search in August 2009: roughly three hundred boxes examined, hundreds of
pages produced, and nothing on Kecksburg.

What the search did produce is an inventory of absence, and this is the part of the case I find
genuinely unsatisfying. Two boxes catalogued 68A2062, covering 1962--67 and described as containing
analysis of fragments to determine national ownership --- the so-called Fragology files, which is
exactly the file class a recovered object would generate --- have been missing since 1987. Box 7 of
accession 70A4099, NASA--DOD material for 1964--66, was destroyed. Seven boxes were signed out in
1996 by a NASA records manager, Paul M.\ Willis, and never returned.

\begin{funfact}
The honest reading of the Kecksburg file loss is not sinister, and it is worse. Federal records
management in the 1980s and 1990s was appalling: boxes were checked out on paper, retention
schedules were applied by clerks with no idea what was in them, and destruction was routine and
logged badly. A conspiracy would have to be competent. What the record shows is an agency that
cannot account for its own filing cabinets, which means it cannot prove innocence either. Loss of
evidence is not evidence --- but an institution that loses evidence forfeits the right to be taken
at its word.
\end{funfact}

\subsection*{Why this is in the same chapter as Die Glocke}

Because the documentaries put it there. The Discovery programme \emph{Nazi UFO Conspiracy} in 2008
and \emph{Ancient Aliens} in 2011 both promoted the claim that the object at Kecksburg \emph{was}
Die Glocke --- captured in 1945, taken to the United States, and either flown or dropped twenty
years later. The chain runs Die Glocke to Project Paperclip (Section~\ref{sec:paperclip}) to
Kecksburg, and it is presented as a single argument.

It is not one. Look at what each link is made of. Die Glocke rests on one journalist's uncorroborated
handwritten notes from documents nobody has seen, as Section~\ref{sec:glocke} sets out. Paperclip is
documented American policy, fully sourced, and what those engineers brought with them was
liquid-fuelled rocketry rather than field propulsion. Kecksburg has real witnesses, a real fireball,
and a real records failure. Joining them produces a story in which the strongest link is asked to
carry the weakest: Paperclip is true, therefore Die Glocke is plausible, therefore the thing on the
flatbed was the Bell. Nothing in that sentence follows from the thing before it. Paperclip proves
that the United States imported German rocket engineers and concealed their war records. It proves
nothing whatever about an object that does not have a single verified document behind it.

Strip the Bell out and Kecksburg is still interesting, which is my point. What remains is a bolide
over six states, a village told to go home, a covered object on a truck consistently described by a
boy at a window across thirteen years, badged men whose agency has no file, a missing radar log, and
two boxes about determining the national ownership of fragments that walked out of a federal archive
in 1987. That is not proof of a craft. It is proof that the United States recovered something it has
never accounted for, and that is a smaller claim I can actually defend.

\begin{srcnotes}
\item Gordon, S.\ Kecksburg case files and USAF documents obtained under FOIA; Air Force response
regarding the Oakdale search team.
\item Chamberlain, V.\ D.\ \& Krause, D.\ J.\ (1967). The Fireball of December 9, 1965.
\emph{Journal of the Royal Astronomical Society of Canada}, 61.
\item NASA National Space Science Data Center, Kosmos 96 entry (re-entry timing and trajectory).
\item Kean, L.\ v.\ NASA, FOIA litigation 2003--2007; NASA records search completed August 2009,
including accession inventories 68A2062 and 70A4099.
\item Oberg, J.\ (2008). Analysis of the Kecksburg recovery accounts and 1960s debris-recovery
practice.
\item \emph{Unsolved Mysteries} (1990) and \emph{The New Roswell: Kecksburg Exposed}, Sci-Fi
Channel (2003), for the Hays, Romansky, Mayes and Weaver testimony.
\end{srcnotes}
